% Part: proof-theory
% Chapter: cut-elimination
% Section: ce-largest
%READERNOTE{SOL6-A294: उलटवण्यात प्रत्येक कोटीच्या कटांची संख्या वाढत नाही, ही सिद्ध सीमा वापरली आहे. अणु-कटातील विगमन कमी उंचीच्या तत्काळ उपसिद्धतेवर लावले आहे. कमी कोटीचे कट असतानाच्या दुर्बलीकरण व आकुंचनाचा आवश्यक आधार स्पष्ट केला आहे; सरावाचे प्रसंग सोडवलेले नाहीत.}

\documentclass[../../../include/open-logic-section]{subfiles}

\begin{document}

\olfileid{pt}{cut}{inv}

\olsection{सर्वोच्च कोटीचे कट काढून टाकणे}

\olref[seq][inv]{lem:G3c-invert} च्या सिद्धतेत दिसले
की \Log{G3c} ने एखाद्या तार्किक नियमाचा निष्कर्ष
(तो नियम \RightR{\lexists} किंवा \LeftR{\lforall}
नसेल तर) !!{prove}s केला, तर त्या नियमाची प्रत्येक
आधार-क्रमवर्तीसुद्धा सिद्धतेची !!{height} न वाढवता
!!{prove}s करता येते. याहून अधिक दाखवता येते:
हेच $\Log{G3c} + \CutCS$ साठीही लागू होते.

आधीप्रमाणे, \CutCS{} अनुमानाची \emph{कट-कोटी}
म्हणजे त्यातील कटच्या !!{formula} ची खोली.

\begin{lem}\ollabel{lem:inv-G3c-cut}
$\Log{G3c} + \CutCS$ मध्ये \RightR{\lexists}
किंवा \LeftR{\lforall} सोडून अन्य तार्किक नियम~$R$
चा निष्कर्ष !!a{proof}~$\pi$ ने !!{prove}s होत असेल,
तर त्या नियमाची प्रत्येक आधार-क्रमवर्ती
!!{proof}~$\pi'$ ने (किंवा नियमाला दोन आधार
असल्यास !!{proof}s~$\pi'$, $\pi''$ यांनी)
!!{prove}s करता येते. शिवाय,
(a)~$\pheight{\pi'}\le\pheight{\pi}$ (आणि
$\pheight{\pi''}\le\pheight{\pi}$);
तसेच (b)~प्रत्येक कोटीसाठी $\pi'$ मधील
(आणि स्वतंत्रपणे $\pi''$ मधील) त्या कोटीच्या
\CutCS{} अनुमानांची संख्या $\pi$ मधील संख्येपेक्षा
मोठी नसते.
\end{lem}

\begin{proof}
\cref{pt:seq:inv:lem:G3c-invert,pt:seq:inv:lem:invert-quant}
यांच्या सिद्धता तपासल्यास हा दावा मिळतो. ती
सिद्धता~$\pheight{\pi}$ वरील विगमनाने होती.
आधार प्रसंगात एक स्वयंसिद्धक दुसऱ्या
स्वयंसिद्धकाने बदलले होते; म्हणून (a) आणि (b)
लगेच मिळतात. $\pi$ चा शेवटचा नियम~$R$
असेल, तर त्याच्या तत्काळ उप-!!{proof}s या
आपल्याला हव्या असलेल्या !!{proof}s~$\pi_i$
आहेत; त्यांची उंची लहान असते आणि प्रत्येक कोटीच्या
कटांची संख्या संपूर्ण $\pi$ मधील संख्येपेक्षा मोठी नसते.
दोन आधारांच्या नियमात एक आधार निवडल्यास दुसऱ्या
शाखेतील कट राहतीलच असे नाही; (b) साठी समानता नको,
ही वरची सीमाच आवश्यक आहे.
इतर प्रसंगांत $\pi$ च्या तत्काळ उप-!!{proof}s वर,
म्हणजे शेवटच्या अनुमान~$R$ च्या आधारांकडे
नेणाऱ्या सिद्धतांवर, विगमन गृहीतक लावून
परिणामांवर तोच नियम $R$ लावला होता.
नव्या सिद्धतेच्या तत्काळ उप-!!{proof}s साठी
(a) व (b) खऱ्या असल्याने संपूर्ण
सिद्धतेसाठीही त्या खऱ्या आहेत. शेवटचे
अनुमान~\CutCS{} असेल तो प्रसंग तपासणे
उरते. पुन्हा फक्त एक उदाहरण पाहू:
$\pi$ ही \LeftR{\land} च्या
$!B \land !C, \Gamma \Sequent \Delta$
या निष्कर्षाची !!a{proof} असून
\CutCS{} ने संपते असे समजा:
\[
\AxiomC{}
\RightLabel{$\pi_1$}
\Deduce$!B \land !C, \Gamma \fCenter \Delta, !A$
\AxiomC{}
\RightLabel{$\pi_2$}
\Deduce$!A, !B \land !C, \Gamma \fCenter \Delta$
\RightLabel{\CutCS}
\BinaryInf$!B \land !C, \Gamma \fCenter \Delta$
\DisplayProof
\]
विगमन गृहीतक $\pi_1$ आणि~$\pi_2$
यांना लागू होते; त्या दोन्ही
$\LeftR{\land}$ च्या निष्कर्षाच्या
उदाहरणांच्या !!{proof}s आहेत.
त्यातून अनुक्रमे $!B, !C, \Gamma \fCenter
\Delta, !A$ आणि $!A, !B, !C, \Gamma \fCenter \Delta$
यांच्या !!{proof}s $\pi_1'$ आणि~$\pi_2'$
मिळतात. \CutCS{} वापरून त्यांना जोडल्यास
पुढील~$\pi'$ मिळते:
\[
\AxiomC{}
\RightLabel{$\pi_1'$}
\Deduce$!B, !C, \Gamma \fCenter \Delta, !A$
\AxiomC{}
\RightLabel{$\pi_2'$}
\Deduce$!A, !B, !C, \Gamma \fCenter \Delta$
\RightLabel{\CutCS}
\BinaryInf$!B, !C, \Gamma \fCenter \Delta$
\DisplayProof
\]
येथे (a) आणि (b) स्पष्टपणे पूर्ण होतात.
\end{proof}

\CutCS{} ऐवजी \Cut{} वापरला असता ही
पद्धत चालली नसती: शेवटच्या !!{proof}
च्या निष्कर्षातील पूर्वांगात
$!B, !C, \Gamma$ च्या दोन प्रती आणि
उत्तरांगात $\Delta$ च्या दोन प्रती आल्या
असत्या. मग आकुंचनाची ग्राह्यता वापरावी
लागली असती; पण
\olref[seq][inv]{prop:G3c-cont-adm}
ची सिद्धता या उलटवण्याच्या उपप्रमेयावरच
आधारित आहे. त्यामुळे युक्तिवाद चक्रीय ठरला असता.
या उलटवण्याच्या उपप्रमेयानंतर दुर्बलीकरण आणि आकुंचनाची
ग्राह्यताही $\Log{G3c}+\CutCS$ पर्यंत वाढवता येते;
त्यामुळे पुढील सिद्धतेत कमी कोटीचे कट असतानाही आकुंचन
वापरता येते. आधीच्या उंचीवरील विगमनात एवढाच नवा प्रसंग
आहे: अंतिम नियम $\CutCS$ असेल, तर दोन्ही आधारांतील
सामायिक संदर्भावर विगमन गृहीतक लावून तोच कट पुन्हा
जोडायचा. कट-सूत्र आणि त्याची कोटी बदलत नाहीत. इतर
नियमांसाठी आधीची सिद्धता वापरायची; आकुंचनाच्या मुख्य
प्रसंगांत आता वरील विस्तारित उलटवणे वापरायचे. त्यातून
उंची आणि कटांच्या सर्वोच्च कोटीत वाढ होत नाही.
संख्यक नियमांच्या आयगेन चरांची नावे आवश्यकतेनुसार आधी
बदलून नव्या संदर्भाच्या चरांपासून वेगळी ठेवायची.

\cref{pt:cut:inv:lem:inv-G3c-cut} वापरून
कट-निर्मूलनाची पर्यायी सिद्धता देता येते.
या सिद्धतेत आपण सर्वोच्च स्थानी असलेली
\CutCS{} अनुमाने एकामागोमाग काढत नाही;
त्याऐवजी सर्वोच्च कोटीची \CutCS{}
अनुमाने काढतो. म्हणजे
$\Log{G3c} + \CutCS$ मधील !!a{proof}~$\pi$
घेऊन तिच्यातील कोणताही कट काढतो
(कदाचित त्याजागी कमी कोटीचे कट येतील)
आणि एकही कट उरेपर्यंत असेच करत राहतो.
अट एवढीच की निवडलेल्या कटच्या वर
त्याच किंवा त्याहून मोठ्या कोटीचे
कट नसावेत.

\begin{lem}\ollabel{lem:max-cut-red-G3c}
$\Log{G3c}+\CutCS$ मधील !!a{proof}~$\pi$
ही $n$ कोटीच्या \CutCS{} अनुमानाने
संपत असेल आणि इतरत्र फक्त
$\Log{G3c}$ चे नियम व $<n$ कोटीची
\CutCS{} अनुमाने वापरत असेल, तर
$\Log{G3c} + \CutCS$ मध्ये
त्याच अंतिम क्रमवर्तीची अशी सिद्धता~$\pi'$
असते की तिच्यातील प्रत्येक \CutCS{}
अनुमानाची कोटी~$<n$ असते.
\end{lem}

\begin{proof}
!!{proof}~$\pi$ चा शेवट पुढीलप्रमाणे असतो:
\[
\AxiomC{}
\RightLabel{$\pi_1$}
\Deduce$\Gamma \fCenter \Delta, !A$
\AxiomC{}
\RightLabel{$\pi_2$}
\Deduce$!A, \Gamma \fCenter \Delta$
\RightLabel{\CutCS}
\BinaryInf$\Gamma \fCenter \Delta$
\DisplayProof
\]
आता~$!A$ च्या रूपानुसार प्रसंग वेगळे करू.

$!A$ हे अणुसूत्र आहे असे समजा. $\pi_1$
आणि $\pi_2$ मधील सर्व \CutCS{} अनुमानांची कोटी
$< \depth{!A} = 0$ असली पाहिजे; म्हणून त्यांपैकी
कोणत्याही सिद्धतेत \CutCS{} अनुमान असू शकत नाही.
म्हणून $\pi_1$ आणि $\pi_2$ दोन्ही कटविरहित आहेत.
$\pheight{\pi_2}$ वरील विगमनाने
$\Gamma \Sequent \Delta$ ची कटविरहित सिद्धता
आहे हे दाखवू. $\pheight{\pi_2} = 0$ असेल,
तर $!A, \Gamma \Sequent \Delta$ हे स्वयंसिद्धक आहे.
$\lfalse$ हे $\Gamma$ मध्ये असेल, तर निष्कर्ष स्वतःच
असत्य-स्वयंसिद्धक आहे. अन्यथा $!A$ मुख्य नसेल, तर
काही अणुसूत्र $!C$ हे $\Gamma$ आणि $\Delta$ या दोन्हींत
!!a{element} असते; म्हणून
$\Gamma \Sequent \Delta$ हे स्वतःच स्वयंसिद्धक
आहे. $!A\ident\lfalse$ असल्याचा प्रसंग
\olref[top]{lem:cut-adm-G3c} च्या सिद्धतेतील प्रसंग B प्रमाणे
हाताळता येतो: कटविरहित $\pi_1$ च्या उत्तरांगातील कटासाठीची
असत्याची प्रत काढल्यास हवी ती सिद्धता मिळते.
उरलेल्या मुख्य अणुसूत्राच्या प्रसंगात
$\Delta = \Delta', !A$. या प्रसंगात $\pi_1$
चा शेवट $\Gamma \Sequent \Delta', !A, !A$
असा असतो. \olref[seq][inv]{prop:G3c-cont-adm},
म्हणजे आकुंचनाची ग्राह्यता, लावल्यास
$\Gamma \Sequent \Delta', !A$ ची
कटविरहित सिद्धता मिळते.
$\pheight{\pi_2} > 0$ असेल, तर
सिद्धता नियम~$R$ ने संपते. त्या नियमाची
एक आधार-क्रमवर्ती घ्या: तिचे रूप
$!A, \Gamma', \Pi \Sequent \Delta', \Lambda$
असते, जेथे $\Pi$ व $\Lambda$ ही
$R$ ची सक्रिय !!{formula}s आहेत.
\olref[seq][adm]{prop:weak-G3c-adm}
हे $\pi_1$ आणि निवडलेल्या आधारावर संपणाऱ्या
$\pi_2$ च्या तत्काळ उपसिद्धतेला लावल्यास
$\Gamma, \Gamma', \Pi \Sequent
\Delta, \Delta', \Lambda, !A$ आणि
$!A, \Gamma, \Gamma', \Pi \Sequent
\Delta, \Delta', \Lambda$ यांच्या
!!{proof}s मिळतात. दुसरी सिद्धता $\pi_2$ पेक्षा
कमी उंचीची राहते; म्हणून त्यांना विगमन गृहीतक लागू होते. त्यामुळे~$R$ च्या प्रत्येक
आधारासाठी $\Gamma, \Gamma', \Pi \Sequent
\Delta, \Delta', \Lambda$ ची !!a{proof}
मिळते. $R$ लावल्यास
$\Gamma, \Gamma \Sequent \Delta, \Delta$
मिळते; आणि \olref[seq][inv]{prop:G3c-cont-adm}
मुळे $\Gamma \Sequent \Delta$ ची
कटविरहित सिद्धता मिळते.
%READERNOTE{OLINC-339: स्रोताचे Delta=Delta,A हे चक्रीय समीकरण; आवश्यक विभागणी Delta=Delta',A अशी दुरुस्त केली आहे.}

$!A$ अणुसूत्र नसेल, तर त्याच्या !!{main
operator} नुसार प्रसंग पाहू. येथे $!A$ हाच
कट-सूत्र आहे. प्रत्येक प्रसंगात
\olref{lem:inv-G3c-cut} वापरून $!A$ मुख्य सूत्र
असलेल्या डाव्या व उजव्या नियमांच्या आधारांच्या
!!{proof}s मिळवू. त्यानंतर
\olref[top]{lem:cut-adm-G3c} च्या सिद्धतेतील
प्रसंग~E प्रमाणे पुढे जाऊ.
\begin{enumerate}
  \item $!A \ident \lnot !B$. !!{proof}s $\pi_1$
  आणि $\pi_2$ यांचे शेवट अनुक्रमे
  $\Gamma \Sequent \Delta, \lnot !B$ आणि
  $\lnot !B, \Gamma \Sequent \Delta$ असे आहेत.
  \olref{lem:inv-G3c-cut} नुसार
  $!B, \Gamma \Sequent \Delta$ आणि
  $\Gamma \Sequent \Delta, !B$ यांच्या
  !!{proof}s $\pi_1'$ आणि $\pi_2'$ मिळतात.
  !!{proof}~$\pi'$ पुढीलप्रमाणे आहे:
  \[
  \AxiomC{}
  \RightLabel{$\pi_2'$}
  \Deduce$\Gamma \fCenter \Delta, !B$
  \AxiomC{}
  \RightLabel{$\pi_1'$}
  \Deduce$!B, \Gamma \fCenter \Delta$
  \RightLabel{\CutCS}
  \BinaryInf$\Gamma \fCenter \Delta$
  \DisplayProof
  \]
  \item $!A \ident (!B \lor !C)$.
  !!{proof}s $\pi_1$ आणि $\pi_2$
  यांचे शेवट
  $\Gamma \Sequent \Delta, !B \lor !C$ आणि
  $!B \lor !C, \Gamma \Sequent \Delta$ असे आहेत.
  $\pi_1$ ला \RightR{\lor} चे उलटवणे
  लावून \olref{lem:inv-G3c-cut} नुसार
  $\Gamma \Sequent \Delta, !B, !C$
  ची !!a{proof} $\pi_1'$ मिळते.
  $\pi_2$ ला \LeftR{\lor} चे उलटवणे
  लावून \olref{lem:inv-G3c-cut} नुसार
  $!B, \Gamma \Sequent \Delta$ ची
  !!a{proof} $\pi_2'$ आणि
  $!C, \Gamma \Sequent \Delta$ ची
  !!a{proof} $\pi_2''$ मिळते.
  \olref[seq][adm]{prop:weak-G3c-adm}
  हे $\pi_2''$ ला लावल्यास
  $!C, \Gamma \Sequent \Delta, !B$
  ची !!a{proof}~$\pi_2'''$ मिळते.
  !!{proof}~$\pi'$ पुढीलप्रमाणे आहे:
  \[
  \AxiomC{}
  \RightLabel{$\pi_1'$}
  \Deduce$\Gamma \fCenter \Delta, !B, !C$
  \AxiomC{}
  \RightLabel{$\pi_2'''$}
  \Deduce$!C, \Gamma \fCenter \Delta, !B$
  \RightLabel{\CutCS}
  \BinaryInf$\Gamma \fCenter \Delta, !B$
  \AxiomC{}
  \RightLabel{$\pi_2'$}
  \Deduce$!B, \Gamma \fCenter \Delta$
  \RightLabel{\CutCS}
  \BinaryInf$\Gamma \fCenter \Delta$
  \DisplayProof
  \]
  \item $!A \ident (!B \land !C)$.
  हा प्रसंग सरावासाठी आहे.
  \item $!A \ident (!B \lif !C)$.
  !!{proof}s $\pi_1$ आणि $\pi_2$
  यांचे शेवट
  $\Gamma \Sequent \Delta, !B \lif !C$ आणि
  $!B \lif !C, \Gamma \Sequent \Delta$ असे आहेत.
  $\pi_1$ ला \RightR{\lif} चे उलटवणे
  लावून \olref{lem:inv-G3c-cut} नुसार
  $!B, \Gamma \Sequent \Delta, !C$
  ची !!a{proof} $\pi_1'$ मिळते.
  $\pi_2$ ला \LeftR{\lif} चे उलटवणे
  लावून \olref{lem:inv-G3c-cut} नुसार
  $\Gamma \Sequent \Delta, !B$ ची
  !!a{proof} $\pi_2'$ आणि
  $!C, \Gamma \Sequent \Delta$ ची
  !!a{proof} $\pi_2''$ मिळते.
  \olref[seq][adm]{prop:weak-G3c-adm}
  हे $\pi_2''$ ला लावल्यास
  $!C, !B, \Gamma \Sequent \Delta$
  ची !!a{proof}~$\pi_2'''$ मिळते.
  !!{proof}~$\pi'$ पुढीलप्रमाणे आहे:
  \[
  \AxiomC{}
  \RightLabel{$\pi_2'$}
  \Deduce$\Gamma \fCenter \Delta, !B$
  \AxiomC{}
  \RightLabel{$\pi_1'$}
  \Deduce$!B, \Gamma \fCenter \Delta, !C$
  \AxiomC{}
  \RightLabel{$\pi_2'''$}
  \Deduce$!C, !B, \Gamma \fCenter \Delta$
  \RightLabel{\CutCS}
  \BinaryInf$!B, \Gamma \fCenter \Delta$
  \RightLabel{\CutCS}
  \BinaryInf$\Gamma \fCenter \Delta$
  \DisplayProof
  \]
  \item $!A \ident \lforall[x][!B(x)]$.
  !!{proof}s $\pi_1$ आणि $\pi_2$
  यांचे शेवट
  $\Gamma \Sequent \Delta, \lforall[x][!B(x)]$
  आणि $\lforall[x][!B(x)], \Gamma \Sequent \Delta$
  असे आहेत. \olref{lem:inv-G3c-cut} नुसार
  $\Gamma \Sequent \Delta, !B(c)$ ची
  !!a{proof}~$\pi_1'$ मिळते.

आता !!{proof}~$\pi_2$ पाहू. तिचा शेवट
$\lforall[x][!B(x)], \Gamma \Sequent \Delta$
असा आहे. $\pi_2$ च्या अंतिम क्रमवर्तीपासून
वर जाताना, डाव्या बाजूला आलेली
$\lforall[x][!B(x)]$ ची अशी प्रत्येक आस्थिती
चिन्हांकित करा की ती अंतिम क्रमवर्तीतील
$\lforall[x][!B(x)]$ च्या आस्थितीपर्यंत
“पोहोचते”. म्हणजे:
(a)~$\pi_2$ च्या अंतिम क्रमवर्तीतील
$\lforall[x][!B(x)]$ चिन्हांकित करा.
(b)~एखाद्या नियमाच्या निष्कर्षाच्या
संदर्भात $\lforall[x][!B(x)]$ चिन्हांकित
असेल, तर त्या नियमाच्या आधारात किंवा
आधारांत तिच्याशी अनुरूप आस्थिती चिन्हांकित
करा. (c)~\LeftR{\lforall} नियमाच्या
निष्कर्षात $\lforall[x][!B(x)]$ हे मुख्य
सूत्र असून चिन्हांकित असेल, तर त्याच्या
आधारातील $\lforall[x][!B(x)]$ आणि
$!B(t)$ यांच्या अनुरूप सक्रिय आस्थिती
चिन्हांकित करा. हा मागोवा $\pi_2$
च्या सिद्धतावृक्षावरच घेतला आहे.

आता $\lforall[x][!B(x)]$ च्या
या सर्व चिन्हांकित आस्थिती पाहू. त्यांपैकी
कोणतीही \LeftR{\lforall} व्यतिरिक्त
अन्य नियमात सक्रिय नसते
(फक्त \LeftR{\lforall} ची सक्रिय
!!{formula}s चिन्हांकित होतात).
$\pi_2$ मधील $\lforall[x][!B(x)]$
च्या सर्व चिन्हांकित आस्थिती काढून टाका.
यातून योग्य !!{proof} मिळाली, तर तिचा शेवट
$\Gamma \Sequent \Delta$ असा आहे आणि
ती $\pi$ च्या जागी ठेवता येते.
याने सर्वोच्च कोटीचा एक कट काढला.

अन्यथा मिळालेली सिद्धता योग्य नाही:
काही \LeftR{\lforall} अनुमानांत सक्रिय
असलेली $\lforall[x][!B(x)]$ ही
!!{formula}s आपण काढून टाकली आहेत.
त्यामुळे पुढील रूपाची “अनुमाने” मिळाली:
\[
  \AxiomC{}
  \RightLabel{$\pi_t$}
  \Deduce$!B(t), \Pi \fCenter \Lambda$
  \RightLabel{\LeftR{\lforall}}
  \UnaryInf$\Pi \fCenter \Lambda$
  \DisplayProof
  \]
अशा प्रत्येक सर्वोच्च “अनुमाना”च्या जागी
पुढील सिद्धता ठेवा:
\[
  \AxiomC{}
  \RightLabel{$\Subst{\pi_1'}{t}{c}[\Pi \Sequent \Lambda]$}
  \Deduce$\Gamma, \Pi \fCenter \Delta, \Lambda, !B(t)$
  \AxiomC{}
  \RightLabel{$\pi_t[\Gamma \Sequent \Delta]$}
  \Deduce$!B(t), \Gamma, \Pi \fCenter \Delta, \Lambda$
  \RightLabel{\CutCS}
  \BinaryInf$\Gamma, \Pi \fCenter \Delta, \Lambda$
  \DisplayProof
  \]
तिच्या खालील प्रत्येक क्रमवर्तीच्या
डावीकडे $\Gamma$ आणि उजवीकडे
$\Delta$ घाला. आधारात चिन्हांकित
$!B(t)$ असलेली सर्व
\LeftR{\lforall} “अनुमाने”
काही~$!B(t)$ वरील \CutCS{}
अनुमानाने बदलेपर्यंत हे पुन्हा करा.
मिळालेल्या !!{proof} ची अंतिम क्रमवर्ती
(आता ती खरोखर योग्य !!{proof} आहे)
$\Gamma, \dots, \Gamma \Sequent
\Delta, \dots, \Delta$ या रूपाची असते.
आकुंचनाची ग्राह्यता लावून तिचे
$\Gamma \Sequent \Delta$ वरील
!!a{proof} मध्ये रूपांतर करा आणि
ती $\pi$ च्या जागी ठेवा. आपण
$\lforall[x][!B(x)]$ वरील एक कट
$!B(t)$ वरील (कदाचित अनेक) कटांनी
बदलला आहे; पण या सर्व कटांची कोटी
कमी आहे. $\pi_1$ किंवा $\pi_2$ मधील उरलेल्या किंवा प्रत करून
वापरलेल्या कटांची कोटीही कमीच आहे. पुढील आकुंचनासाठी
वर सिद्ध केलेली, सर्वोच्च कट-कोटी न वाढवणारी विस्तारित
ग्राह्यता वापरली आहे.
\end{enumerate}
\end{proof}

\begin{prob}
\olref{lem:max-cut-red-G3c} च्या सिद्धतेतील
$!A \ident (!B \land !C)$ हा प्रसंग तपासा.
म्हणजे, दोन उपसिद्धतांसह पुढीलप्रमाणे संपणारी
!!a{proof}~$\pi$ असेल:
\[
\AxiomC{}
\RightLabel{$\pi_1$}
\Deduce$\Gamma \fCenter \Delta, !B \land !C$
\AxiomC{}
\RightLabel{$\pi_2$}
\Deduce$!B \land !C, \Gamma \fCenter \Delta$
\RightLabel{\CutCS}
\BinaryInf$\Gamma \fCenter \Delta$
\DisplayProof
\]
आणि $\pi_1$ व $\pi_2$ यांतील
सर्व कटांची कोटी $<\depth{!B
\land !C}$ असेल, तर
$\Gamma \Sequent \Delta$ ची
फक्त $<\depth{!B \land !C}$
कोटीचे कट असलेली !!a{proof}~$\pi'$
असते हे दाखवा.
\end{prob}

% \begin{prob}
% \olref{lem:max-cut-red-G3c} च्या सिद्धतेतील
% $!A \ident \lexists[x][!B(x)]$ हा प्रसंग तपासा.
% \end{prob}

\olref{lem:max-cut-red-G3c} नुसार
!!a{proof} मधील सर्वोच्च कोटीचे
\CutCS{} अनुमान कमी कोटीच्या
\CutCS{} अनुमानांनी बदलता येते.
हे उपप्रमेय पुन्हा वापरून,
सर्वोच्च~\CutCS{} ने संपणाऱ्या
सर्वोच्च उप-!!{proof}s वर ते एकामागोमाग
लावल्यास कोणत्याही !!a{proof} मधून कितीही
\CutCS{} अनुमाने काढता येतात.
%READERNOTE{OLINC-340: स्रोत येथे सर्वोच्च कटासाठीचे वेगळे उपप्रमेय दाखवतो; या परिच्छेदातील कोटी-कमीकरण मागील उपप्रमेयावर आधारित आहे.}

\begin{thm}\ollabel{thm:cut-elim-G3c-largest}
$\Log{G3c} + \CutCS$ मधील कोणतीही
!!{proof} $\pi$ ही~\Log{G3c} मधील
सिद्धतेत रूपांतरित करता येते.
\end{thm}

\begin{proof}
  प्रथम $\pi$ मधील सर्वोच्च कोटीचे
  सर्व कट काढता येतात हे सिद्ध करू.
  $d$ ही $\pi$ मधील कटांची सर्वोच्च
  कोटी आणि $n$ ही $d$ कोटीच्या कटांची
  संख्या असू द्या. $n$ वरील विगमनाने
  सिद्धता करू. $n=0$ असेल, तर
  सिद्ध करायचे काही नाही. $n>0$
  असेल, तर $d$ कोटीचा सर्वोच्च कट
  निवडा आणि त्यावर संपणारी उप-!!{proof}
  $\pi_1$ असू द्या.
  \olref{lem:max-cut-red-G3c} नुसार
  त्याच अंतिम क्रमवर्तीची, फक्त
  $<d$ कोटीचे कट असलेली,
  !!a{proof}~$\pi_1'$ आहे.
  $\pi$ मध्ये $\pi_1$ च्या जागी
  $\pi_1'$ ठेवा. मिळालेल्या
  !!a{proof} मध्ये $d$ कोटीचे
  $n-1$ कट असतात; म्हणून
  विगमन गृहीतक लागू होते.

आता $d$ वरील विगमनाने निष्कर्ष
मिळतो. $d=0$ असेल, तर सर्व कट
अणुसूत्रांवरील आहेत आणि मागील
निष्कर्षाने ते सर्व काढता येतात.
$d>0$ असेल, तर मागील निष्कर्षाने
$d$ कोटीचे सर्व कट $<d$ कोटीच्या
कटांनी बदललेली सिद्धता मिळते.
$\pi$ मधील कटांची सर्वोच्च कोटी
$d$ असल्यामुळे नव्या सिद्धतेतील
कटांची सर्वोच्च कोटी $<d$ होते;
म्हणून विगमन गृहीतक लागू होते.
\end{proof}

\end{document}
